\documentclass{article}
\usepackage[utf8]{inputenc}
\usepackage{titling}
\usepackage{graphicx}
\usepackage{enumitem}
\usepackage{array}
\usepackage[includeheadfoot, margin=1em,headheight=2em]{geometry}
\usepackage[hidelinks]{hyperref}
\geometry{a4paper, left=2cm, right=2cm, top=2cm, bottom=2cm}
\newcolumntype{P}[1]{>{\centering\arraybackslash}p{#1}}
\renewcommand{\arraystretch}{1.5} % Default value: 1
\setlength{\droptitle}{-6em}

%font
\usepackage[defaultfam,tabular,lining]{montserrat}
\usepackage[T1]{fontenc}
\renewcommand*\oldstylenums[1]{{\fontfamily{Montserrat-TOsF}\selectfont #1}}

%custom bold 
\usepackage[outline]{contour}
\usepackage{xcolor}
\newcommand{\custombold}{\contour{black}}

%table colors
\usepackage{color, colortbl}
\definecolor{Blue}{rgb}{0.51,0.68,0.79}
\definecolor{LightBlue}{rgb}{0.82,0.87,0.90}
\definecolor{LighterBlue}{rgb}{0.93,0.95,0.96}

%Header
\usepackage{fancyhdr, xcolor}
\pagestyle{fancy}
\let\oldheadrule\headrule
\renewcommand{\headrule}{\color{Blue}\oldheadrule}
\renewcommand{\headrulewidth}{0.2em}
\fancyhead[L]{CyberSorcerers - Verbale di riunione interna}
\fancyhead[C]{}
\fancyhead[R]{24/10/2023}

\title{\Huge{\textbf{Verbale di riunione}}\vspace{-1em}}
\date{}

\begin{document}
\maketitle
\vspace{-3em}
\begin{figure}[h]
  \centering
  \includegraphics[width=6cm, height=6cm]{documenti/logo rotondo}
  \label{fig:immagine}
\end{figure}

\begin{center}
\Large{27 Novembre 2023 - 16:00\\
Google Meet\\
Verbale redatto da: Giulia Dentone\\}
\end{center}
\vspace{2em}
\large{
\begin{center}
    \begin{tabular}{P{24em}}
        \rowcolor{Blue}
        \textbf{Partecipanti alla riunione}\\
        \rowcolor{LightBlue}
        \custombold{Giulia Dentone}: presente\\
        \rowcolor{LighterBlue}
        \custombold{Giovanni Moretti}: presente\\
        \rowcolor{LightBlue}
        \custombold{Sabrina Caniato}: presente\\
        \rowcolor{LighterBlue}
        \custombold{Nicola Lazzarin}: presente\\
        \rowcolor{LightBlue}
        \custombold{Samuele Vignotto}: presente \\
        \rowcolor{LighterBlue}
        \custombold{Andrea Rezzi}: presente \\
    \end{tabular}
\end{center}}
\newpage
\section{Obiettivo della riunione}
Suddivisione dei compiti per la partecipazione al PoC.

\section{Agenda}
\label{sec:decisioni}
\subsection{Diario di bordo}
Abbiamo assegnato discusso sul contenuto e la realizzazione del prossimo Diario di Bordo che si terrà in data 29/11/2023.

\section{Decisioni}
\subsection{Jira}
Abbiamo deciso di utilizzare come software di project management Jira.

\subsection{Stand Up Meetings}
Abbiamo deciso, al fine di migliorare il nostro Way of Working, di tenere degli Stand up meetinng interni della durata di massimo 10 minuti due volte a settimana, per aggiornarci sull'avanzamento dei task assegnati durante le riunioni interne. Le date dei meeting saranno indicativamente il mercoledì e il venerdì alle ore 13:00.

\subsection{Analisi dei requisiti}
Abbiamo deciso di suddividere la stesura dell'analisi dei requisiti in to do piccoli in modo da rendere il lavoro più agile.

\subsection{Azienda}
Abbiamo deciso di contattare via mail l'azienda al fine di aprire un canale slack, di avere le credenziali necessarie per iniziare la fase di testing su Chat GPT e Amazon Web Services. Abbiamo inoltre deciso di richiedere una riunione esterna dopo il giorno 8/12/2023, data fissata come limite per preparare il materiale utile al meeting e avere delle domande più attinenti e puntuali da porre all'azienda.

\section{Azioni}
\begin{center}
\begin{tabular}{P{20em} P{10em} P{10em}} 
    \custombold{Task} & \custombold{Responsabile} & \custombold{Termine massimo}\\
  \hline
  Stesura del verbale interno del 27/11/2023 & Giulia Dentone & 28/11/2023 \\ 
  \hline
  Preparazione del diario di bordo del 29/11/2023 & Samuele Vignotto & 28/11/2023 \\ 
  \hline
  Scrivere la mail all'azienda & Nicola Lazzarin & 27/11/2023\\ 
  \hline
  Preparazione per il PoC & Giulia Dentone e Andrea Rezzi & 08/12/2023 \\ 
  \hline
  Procedere con la stesura dell'Analisi dei requisiti & Nicola Lazzarin e Samuele Vignotto & 08/11/2023 \\
  \hline
  Definire la struttura e i potenziali contenuti del Piano di Progetto & Giovanni Moretti & 02/12/2023 \\ 
  \hline
  Definire la stuttura e i potenziali contenuti del Piano di Qualifica & Sabrina Caniato & 02/12/2023 \\ 
  \hline
  Digitalizzare il workflow ideale formitoci dall'azienda& Giulia Dentone & 29/11/2023 \\ 
  \hline
 
  
\end{tabular}
\end{center}
\section{Attività future}
\subsection{Riunione interna:}
La prossima riunione si terrà su Google meet il giorno 04-12-2023 alle ore 15:30.
\section{Eventuali note}
In attesa della risposta dell'azienda per fissare una riunione esterna.
\end{document}